% GENERATED FILE. Edit canonical lessons, then run npm run generate:books.
% canonical-source-hash: fnv1a64:85567a2595eb1ea4
% canonical-lessons: KA-C275-carj-madu, KA-C275-parisilisu, KA-C275-anuvadisu, KA-C275-arji-haku, KA-C275-holisu

\chapter{To charge, To check, To translate, To apply, To compare}
\label{ch:ka-to-charge-to-check-to-translate-to-apply-to-compare}
\hlchaptermodality{275}

\begin{chapteropening}
\textbf{By the end of this chapter:} \emph{I can say to charge, to check, to translate, to apply and to compare.}

Complete the last lesson of chapter 275: I can say to charge, to check, to translate, to apply and to compare.
\end{chapteropening}

\section[\texorpdfstring{\kn{ಚಾರ್ಜ್} \kn{ಮಾಡು} (cārj māḍu)}{ಚಾರ್ಜ್ ಮಾಡು (cārj māḍu)}]{\kn{ಚಾರ್ಜ್} \kn{ಮಾಡು} (cārj māḍu) — to charge (a phone)}

\label{lesson:KA-C275-carj-madu}



\begin{quote}
Before the new one: say the Kannada for to withdraw, then the Kannada for to sign.
\end{quote}

\subsection*{You'll want to know: \kn{ಚಾರ್ಜ್} \kn{ಮಾಡು}}
\textbf{\kn{ಚಾರ್ಜ್} \kn{ಮಾಡು}} — \emph{cārj māḍu} — \textquotedblleft{}to charge (a phone)\textquotedblright{}.

\begin{grammarlens}[title={What you've built}]
One more word for this chapter.
\end{grammarlens}

\subsection*{Your turn}
\begin{itemize}
\raggedright
  \item \emph{Say it:} \emph{cārj māḍu}
  \item \emph{Say it:} \emph{cārj māḍu}, once more
  \item \emph{Say it:} say \emph{cārj māḍu}
  \item \emph{Recall:} say the Kannada for to withdraw, then the Kannada for to sign, then say \emph{cārj māḍu} again
\end{itemize}

\begin{tcolorbox}[breakable,colback=teal!4,colframe=teal!35!black,title={Before you move on}]
What does \kn{ಚಾರ್ಜ್} \kn{ಮಾಡು} mean? (\textquotedblleft{}To charge (a phone)\textquotedblright{}.) Say it once more.
\end{tcolorbox}
\section[\texorpdfstring{\kn{ಪರಿಶೀಲಿಸು} (pariśīlisu)}{ಪರಿಶೀಲಿಸು (pariśīlisu)}]{\kn{ಪರಿಶೀಲಿಸು} (pariśīlisu) — to check, to verify}

\label{lesson:KA-C275-parisilisu}



\begin{quote}
Before the new one: say the Kannada for to sign, then the Kannada for to charge.
\end{quote}

\subsection*{You'll want to know: \kn{ಪರಿಶೀಲಿಸು}}
\textbf{\kn{ಪರಿಶೀಲಿಸು}} — \emph{pariśīlisu} — \textquotedblleft{}to check, to verify\textquotedblright{}.

\begin{grammarlens}[title={What you've built}]
One more word for this chapter.
\end{grammarlens}

\subsection*{Your turn}
\begin{itemize}
\raggedright
  \item \emph{Say it:} \emph{pariśīlisu}
  \item \emph{Say it:} \emph{pariśīlisu}, once more
  \item \emph{Say it:} say \emph{pariśīlisu}
  \item \emph{Recall:} say the Kannada for to sign, then the Kannada for to charge, then say \emph{pariśīlisu} again
\end{itemize}

\begin{tcolorbox}[breakable,colback=teal!4,colframe=teal!35!black,title={Before you move on}]
What does \kn{ಪರಿಶೀಲಿಸು} mean? (\textquotedblleft{}To check, to verify\textquotedblright{}.) Say it once more.
\end{tcolorbox}
\section[\texorpdfstring{\kn{ಅನುವಾದಿಸು} (anuvādisu)}{ಅನುವಾದಿಸು (anuvādisu)}]{\kn{ಅನುವಾದಿಸು} (anuvādisu) — to translate}

\label{lesson:KA-C275-anuvadisu}



\begin{quote}
Before the new one: say the Kannada for to charge, then the Kannada for to check.
\end{quote}

\subsection*{You'll want to know: \kn{ಅನುವಾದಿಸು}}
\textbf{\kn{ಅನುವಾದಿಸು}} — \emph{anuvādisu} — \textquotedblleft{}to translate\textquotedblright{}.

\begin{grammarlens}[title={What you've built}]
One more word for this chapter.
\end{grammarlens}

\subsection*{Your turn}
\begin{itemize}
\raggedright
  \item \emph{Say it:} \emph{anuvādisu}
  \item \emph{Say it:} \emph{anuvādisu}, once more
  \item \emph{Say it:} say \emph{anuvādisu}
  \item \emph{Recall:} say the Kannada for to charge, then the Kannada for to check, then say \emph{anuvādisu} again
\end{itemize}

\begin{tcolorbox}[breakable,colback=teal!4,colframe=teal!35!black,title={Before you move on}]
What does \kn{ಅನುವಾದಿಸು} mean? (\textquotedblleft{}To translate\textquotedblright{}.) Say it once more.
\end{tcolorbox}
\section[\texorpdfstring{\kn{ಅರ್ಜಿ} \kn{ಹಾಕು} (arji hāku)}{ಅರ್ಜಿ ಹಾಕು (arji hāku)}]{\kn{ಅರ್ಜಿ} \kn{ಹಾಕು} (arji hāku) — to apply (for a job, a loan)}

\label{lesson:KA-C275-arji-haku}



\begin{quote}
Before the new one: say the Kannada for to check, then the Kannada for to translate.
\end{quote}

\subsection*{You'll want to know: \kn{ಅರ್ಜಿ} \kn{ಹಾಕು}}
\textbf{\kn{ಅರ್ಜಿ} \kn{ಹಾಕು}} — \emph{arji hāku} — \textquotedblleft{}to apply (for a job, a loan)\textquotedblright{}.

\begin{grammarlens}[title={What you've built}]
One more word for this chapter.
\end{grammarlens}

\subsection*{Your turn}
\begin{itemize}
\raggedright
  \item \emph{Say it:} \emph{arji hāku}
  \item \emph{Say it:} \emph{arji hāku}, once more
  \item \emph{Say it:} say \emph{arji hāku}
  \item \emph{Recall:} say the Kannada for to check, then the Kannada for to translate, then say \emph{arji hāku} again
\end{itemize}

\begin{tcolorbox}[breakable,colback=teal!4,colframe=teal!35!black,title={Before you move on}]
What does \kn{ಅರ್ಜಿ} \kn{ಹಾಕು} mean? (\textquotedblleft{}To apply (for a job, a loan)\textquotedblright{}.) Say it once more.
\end{tcolorbox}
\section[\texorpdfstring{\kn{ಹೋಲಿಸು} (hōlisu)}{ಹೋಲಿಸು (hōlisu)}]{\kn{ಹೋಲಿಸು} (hōlisu) — to compare}

\label{lesson:KA-C275-holisu}



\begin{quote}
Before the new one: say the Kannada for to translate, then the Kannada for to apply.
\end{quote}

\subsection*{You'll want to know: \kn{ಹೋಲಿಸು}}
\textbf{\kn{ಹೋಲಿಸು}} — \emph{hōlisu} — \textquotedblleft{}to compare\textquotedblright{}.

\begin{grammarlens}[title={What you've built}]
One more word for this chapter.
\end{grammarlens}

\subsection*{Your turn}
\begin{itemize}
\raggedright
  \item \emph{Say it:} \emph{hōlisu}
  \item \emph{Say it:} \emph{hōlisu}, once more
  \item \emph{Say it:} say \emph{hōlisu}
  \item \emph{Recall:} say the Kannada for to translate, then the Kannada for to apply, then say \emph{hōlisu} again
\end{itemize}

\begin{tcolorbox}[breakable,colback=teal!4,colframe=teal!35!black,title={Before you move on}]
What does \kn{ಹೋಲಿಸು} mean? (\textquotedblleft{}To compare\textquotedblright{}.) Say it once more.
\end{tcolorbox}
